% Discussion, open problems, acknowledgements, code and data (front writer). Plan: ../OUTLINE.md "§13 Discussion";
% synthesis research/notes-final.md §0 (philosophy paragraph), §2.7 (verdict), §2.8, §6 (open problems; the five starred ones
% are named first), §8.6 (re-runs, for the Code and data statement); CLAIMS.md §12 (every q: label is cited below).
% Acknowledgements as in the companion papers (DECISIONS §A4, REFRAME §4); the body uses no work of S. Eisenstat or
% D. Vaintrob, so they are not thanked here. Author decision (2026-10-17): the sentence "Hänni takes responsibility for the content of this paper" is removed now that the paper has two authors; do not reinstate it.
% Edit pass of 2026-10-17 (review/DECISIONS.md; FRO-13 to FRO-16, FRO-19, FRO-L2): U1 (long form, first sentence shortened as
% allowed), U2 (discussion version), U3 (first sentence); FRO-14, FRO-15; acknowledgements without the sentence that repeated the
% title-page footnote; list A and list B of DECISIONS §L3 (front, §13). The Code and data sentence on the README stays (FRO-19):
% the README is a blocking item of the publication step (DECISIONS §X9).

\section{Discussion}\label{sec:discussion}

\paragraph{What the formalism shows.} \SAI makes precise the view that theories answer questions posed in a given language, and the view survives in comparative form: it changes the target, the data type, logical coherence and partiality (\cref{sec:phil:questions,tab:phil:differences}), while the inference principle is unchanged: version by version and up to constant factors, \SAIabbr agrees with operator induction over consistent labellers (\cref{sec:soi}). On data labelled true or false the idealised formalism cannot see the difference between a theory and a program. In the settings we study, shape becomes visible only when something other than program length is charged (the written theory, the time to produce it, or the derivations) or the class of hypotheses is restricted (to finite theories); there it is relative to a reasoner and a vocabulary, and for theoretical terms what pays is relations and entities, not vocabulary as such. Whether this holds beyond these settings is open (\cref{q:shape:beyond}). That textbook theories are declarative is an observation we accept, and we find it plausible that this shape is good, but in our models declarative against procedural form buys at most a constant factor once intermediate concepts are available and the reasoner can chain through them; what is missing is a formal setting in which that shape provably pays. Finally, ``Kepler's laws are only approximately true'' is a limitation of proof-requiring compatibility with deterministic theories, not of the axiom view as \SAIabbr formalises it: its mixture contains randomised theories, and under its pricing an approximately true law is a law with a calibrated error model (\cref{prop:laws:errormodel}).

\paragraph{What it cannot show.} The axiom view models the declarative core of scientific theories (contents, inferential relations, partiality, approximate truth) and misses their procedural, statistical and dynamic sides (\cref{sec:limits}). It supports no general realist conclusion about the theories it favours: it identifies a complete r.e.\ theory that the data present (\cref{prop:basic:limits}), but its limit posterior can live on random theories and, with adversarially chosen true data, on false ones (\cref{thm:logic:lottery,cor:logic:false}). It evaluates theories and does not find them, and it is not computable.

\subsection{Open problems}\label{sec:discussion:open}

\begin{enumerate}[leftmargin=*,label=(\arabic*)]
\item Does $\MA$ dominate every lower semicomputable provability function of a random consistent theory (\cref{q:basic:domination})?
\item Is there a stream with bounded one-tape loss on which the coin-free loss diverges computably fast, and does bounded one-tape loss on a co-r.e.\ stream imply that a consistent r.e.\ theory proves it (\cref{q:logic:computablefast})?
\item Over a true background containing $\PA$, is the one-tape posterior probability of $\Sigma_1$-soundness bounded away from $1$ uniformly over all data (\cref{q:logic:soundness})?
\item Does the dependence of written size on the pair (reasoner, vocabulary) extend beyond Boolean laws, to real-valued laws with numerical solvers, probabilistic programs with conditioning and simulation codes (\cref{q:shape:beyond})?
\item Is there a formal setting for communities of bounded reasoners who communicate, check and extend theories piece by piece, in which theories of declarative shape provably do better than programs (\cref{q:phil:declarative})?
\end{enumerate}
The further questions are \cref{q:basic:tape,q:basic:indeporder,q:soi:lsc,q:soi:pathwise,q:soi:presburger,q:logic:soundtrue,q:logic:codingO1,q:logic:thn,q:logic:kreisel,q:terms:binaryobs,q:terms:a1,q:terms:unaryfunction,q:shape:minimum,q:shape:learning,q:shape:cellprobe,q:laws:JR,q:laws:domains}.

\paragraph{Directions.} The directions we find most worth pursuing are step 4 of the ladder, for which we have no formal handle; the setting of problem (5), where the declarative shape of written theories may find its explanation; a costed-deduction mixture over unknown axiom systems, combining \SAIabbr, uncertain about the axioms, with logical induction \citep{garrabrant2016logical}, uncertain about their consequences, as induction in mathematics needs (\cref{prop:logic:oldevidence}); and a native treatment of statistical confrontation inside the proof-requiring core (\cref{sec:limits:stat}). Smaller ones: whether explanatory asymmetry is interventional in general, and a derivation of Kepler's laws within stated tolerances from Newton's theory in a prover for validated numerics (\cref{sec:limits:textbook}).

\section*{Acknowledgements}\addcontentsline{toc}{section}{Acknowledgements}
Several of the questions and examples studied here first appeared in H\"anni's working notes \citep{hanni2026notes}. The results of our earlier report \citep{claude2026axiominduction} and of the companion papers \citep{claudeSFI2026fundamentals,claudeSFI2026steeper} that we use were prepared and checked in the same way as this paper. No human has yet independently checked all of the mathematics, and nothing was verified in a proof assistant; readers should weigh the results accordingly.

\section*{Code and data}\addcontentsline{toc}{section}{Code and data}
All computations reported here were done with seeded or deterministic Python scripts (Python~3.11 with NumPy, SciPy and mpmath; the planetary computations use pyerfa, the Python interface to the ERFA library). The scripts, their saved outputs and independent re-implementations of the main computations are in the folder \texttt{solomonoff-axiom-induction/} of the repository \url{https://github.com/kaarelh/AI-works}; its README lists the script behind each reported number and how to re-run it. Re-running the scripts reproduces the saved outputs, apart from lines that report run times; the longest computations, on the Mars pool, take more than half an hour each, and their figures were also reproduced by independent re-implementations.
